\chapter{Detailed Worked Derivations And Examples}
\label{chap:detailed-worked-derivations}

\section{Purpose Of This Chapter}
\label{sec:detailed-purpose}

The earlier chapters develop the structure of mechanics. This chapter slows
down and works through representative calculations in full. The goal is not to
collect tricks. The goal is to show how assumptions, notation, conserved
quantities, reductions, and stability criteria fit together in examples that
students can use as templates.

Each example follows the same pattern:
\[
\begin{gathered}
  \text{model assumptions}
  \longrightarrow
  \text{coordinates and Lagrangian/Hamiltonian}\\
  \longrightarrow
  \text{conserved quantities or constraints}
  \longrightarrow
  \text{reduced equation}
  \longrightarrow
  \text{interpretation}.
\end{gathered}
\]
This pattern is more important than any single example. The most common source
of confusion in advanced mechanics is mixing levels: treating a coordinate
artifact as a physical effect, treating a conserved quantity as an independent
coordinate before checking regularity, or treating a numerical picture as a
theorem. The derivations in this chapter slow down exactly at those points.

For that reason, some calculations repeat ideas that appeared earlier in a more
compressed form. The repetition is intentional. A fast derivation is good for
orientation; a slow derivation is where one sees which assumptions are doing
work, which signs come from conventions, and which quantities are actually
invariant.

\paragraph{How to compare these examples.}
Read each example as a template with replaceable parts. The central-force
example replaces a rotational degree of freedom by angular momentum. The
pendulum replaces an orbit by a phase portrait with a separatrix. The rigid top
replaces a configuration-space picture by motion on invariant tori. The
resonance and Henon-Heiles examples show how small denominators enter. The
fluid and gauge examples show how constraints produce pressure or a mechanical
connection. The calculations are different, but the logic is the same: identify
the reduced variables, derive the diagnostic equation, and interpret the
geometry of the result.

\section{Central Forces From First Principles}
\label{sec:detailed-central-forces}

The central-force problem is the prototype of reduction by symmetry. Rotational
invariance gives conservation of angular momentum; angular momentum confines
the motion to a plane; polar coordinates in that plane reduce the dynamics to a
one-dimensional radial equation with an effective potential. The later Kepler
and precession calculations are refinements of this same reduction.

The lesson is broader than celestial mechanics. Every reduction in these notes
has the same shape: use a symmetry to find a conserved quantity, use the
conserved quantity to remove a variable, and then interpret the remaining
equation on the reduced space. Central forces are the cleanest place to see the
pattern because the geometry is visible in ordinary Euclidean space.

Let a particle of mass \(m\) move in a central potential \(V(r)\), where
\[
  r=\norm{x},
  \qquad
  x\in\R^3\setminus\{0\}.
\]
The Lagrangian is
\[
  L=\frac12m\norm{\dot x}^2-V(r).
\]
Rotational invariance gives angular momentum
\[
  J=x\times p=m x\times\dot x.
\]
Since
\[
  \frac{\dd J}{\dd t}
  =
  \dot x\times m\dot x+x\times m\ddot x
  =
  x\times\left(-V'(r)\frac{x}{r}\right)=0,
\]
the vector \(J\) is constant. Because \(J\cdot x=0\) and
\(J\cdot\dot x=0\), the motion lies in the plane perpendicular to \(J\). Use
polar coordinates \((r,\theta)\) in this plane:
\[
  L=\frac12m(\dot r^2+r^2\dot\theta^2)-V(r).
\]
The angle \(\theta\) is cyclic, so
\[
  \ell=\frac{\partial L}{\partial\dot\theta}=mr^2\dot\theta
\]
is conserved. The energy is
\[
  E=\frac12m\dot r^2+\frac{\ell^2}{2mr^2}+V(r).
\]
Thus the radial motion is one-dimensional motion in
\[
  V_{\mathrm{eff}}(r)
  =
  V(r)+\frac{\ell^2}{2mr^2}.
\]
Here \(\ell\) is the massful angular momentum. To compare with the
specific-energy Kepler convention used in
Section~\ref{sec:kepler-perturbation}, divide the Hamiltonian, energy, and
angular momentum by the particle mass; then
\(\ell_{\mathrm{spec}}=\ell/m=r^2\dot\theta\). The orbit equation is unchanged
after this bookkeeping conversion, provided \(k/m\) is read as the
gravitational parameter \(GM\).

\subsection{Binet Equation}
\label{subsec:detailed-binet}

Let
\[
  u(\theta)=\frac1r.
\]
Since
\[
  \dot\theta=\frac{\ell}{mr^2}=\frac{\ell}{m}u^2,
\]
we have
\[
  \frac{\dd}{\dd t}
  =
  \frac{\ell}{m}u^2\frac{\dd}{\dd\theta}.
\]
Also
\[
  \dot r
  =
  \frac{\dd}{\dd t}(u^{-1})
  =
  -u^{-2}\dot u
  =
  -\frac{\ell}{m}u',
\]
where a prime denotes \(\dd/\dd\theta\). Differentiating once more,
\[
  \ddot r
  =
  -\frac{\ell}{m}\frac{\dd u'}{\dd t}
  =
  -\frac{\ell^2}{m^2}u^2u''.
\]
The radial equation is
\[
  m(\ddot r-r\dot\theta^2)=F(r),
  \qquad
  F(r)=-V'(r).
\]
Substitute the formulas above:
\[
  m\left(
    -\frac{\ell^2}{m^2}u^2u''
    -
    \frac1u\frac{\ell^2}{m^2}u^4
  \right)
  =
  -\frac{\ell^2}{m}u^2(u''+u)
  =
  F(1/u).
\]
Therefore
\[
  \boxed{
  u''+u
  =
  -\frac{m}{\ell^2u^2}F(1/u).
  }
\]
This is Binet's equation. It turns a central-force problem into a differential
equation for the orbit shape \(r(\theta)\).

\subsection{Kepler Conics And The Laplace-Runge-Lenz Vector}
\label{subsec:detailed-kepler-lrl}

For the Kepler potential
\[
  V(r)=-\frac{k}{r},
  \qquad
  k>0,
\]
the force is
\[
  F(r)=-\frac{k}{r^2}.
\]
Binet's equation becomes
\[
  u''+u=\frac{mk}{\ell^2}.
\]
The solution is
\[
  u(\theta)=\frac{mk}{\ell^2}\left[1+e\cos(\theta-\theta_0)\right],
\]
or
\[
  r(\theta)=\frac{\ell^2/(mk)}
  {1+e\cos(\theta-\theta_0)}.
\]
This is a conic section with focus at the force center.

The extra conserved vector is
\[
  A=p\times J-mk\frac{x}{r}.
\]
To verify conservation, write \(\hat r=x/r\). Since \(J\) is constant,
\[
  \dot A=\dot p\times J-mk\dot{\hat r}.
\]
The force law gives
\[
  \dot p=-\frac{k}{r^2}\hat r.
\]
Use
\[
  J=mx\times\dot x=mr^2\dot\theta\,\hat z=\ell\hat z
\]
in the orbital plane. Then
\[
  \dot{\hat r}=\dot\theta\,\hat\theta,
\]
and
\[
  \dot p\times J
  =
  -\frac{k}{r^2}\hat r\times(\ell\hat z)
  =
  \frac{k\ell}{r^2}\hat\theta
  =
  mk\dot\theta\,\hat\theta
  =
  mk\dot{\hat r}.
\]
Hence
\[
  \dot A=0.
\]
The vector \(A\) points toward periapse. Taking the dot product with \(x\),
\[
  A\cdot x
  =
  (p\times J)\cdot x-mkr.
\]
By cyclic permutation,
\[
  (p\times J)\cdot x=J\cdot(x\times p)=J\cdot J=\ell^2.
\]
Therefore
\[
  A r\cos(\theta-\theta_0)=\ell^2-mkr,
\]
which rearranges to
\[
  r=\frac{\ell^2/(mk)}{1+(A/(mk))\cos(\theta-\theta_0)}.
\]
Thus the eccentricity is
\[
  e=\frac{\norm A}{mk}.
\]
The conic equation and the hidden conservation law are the same fact written
in two languages.

\subsection[Small Periapse Precession]{Small Periapse Precession From A
\texorpdfstring{\(1/r^3\)}{1/r^3} Perturbation}
\label{subsec:detailed-precession}

A useful perturbation of the Kepler potential is
\[
  V(r)=-\frac{k}{r}-\frac{\epsilon}{r^3},
  \qquad
  0<|\epsilon|\ll1.
\]
The radial force is
\[
  F(r)=-V'(r)=-\frac{k}{r^2}-\frac{3\epsilon}{r^4}.
\]
Binet's equation gives
\[
  u''+u
  =
  \frac{mk}{\ell^2}
  +
  \frac{3m\epsilon}{\ell^2}u^2.
\]
For a nearly circular orbit, write
\[
  u(\theta)=u_0+\eta(\theta),
  \qquad
  |\eta|\ll u_0.
\]
The constant \(u_0\) solves
\[
  u_0=\frac{mk}{\ell^2}+\frac{3m\epsilon}{\ell^2}u_0^2.
\]
Linearizing in \(\eta\),
\[
  \eta''+\eta
  =
  \frac{6m\epsilon u_0}{\ell^2}\eta,
\]
or
\[
  \eta''+\kappa^2\eta=0,
  \qquad
  \kappa^2=1-\frac{6m\epsilon u_0}{\ell^2}.
\]
The radial oscillation has angular period
\[
  \Delta\theta_r=\frac{2\pi}{\kappa}.
\]
The periapse advance per radial cycle is the excess over \(2\pi\):
\[
  \Delta\varpi
  =
  2\pi\left(\frac1\kappa-1\right)
  \approx
  \frac{6\pi m\epsilon u_0}{\ell^2}.
\]
Thus the \(1/r^3\) correction breaks the closed Kepler ellipse and produces
precession. This is the simplest analytic counterpart to the numerical
precession demo.

\section{Pendulum Phase Space And The Separatrix}
\label{sec:detailed-pendulum}

The simple pendulum is the canonical one-degree-of-freedom Hamiltonian example.
Let \(\theta\) be the angle from the downward vertical. The Lagrangian is
\[
  L=\frac12m\ell^2\dot\theta^2+mg\ell\cos\theta.
\]
The momentum and Hamiltonian are
\[
  p_\theta=m\ell^2\dot\theta,
  \qquad
  H(\theta,p_\theta)
  =
  \frac{p_\theta^2}{2m\ell^2}-mg\ell\cos\theta.
\]
Hamilton's equations are
\[
  \dot\theta=\frac{p_\theta}{m\ell^2},
  \qquad
  \dot p_\theta=-mg\ell\sin\theta.
\]
The equilibria are
\[
  (\theta,p_\theta)=(0,0)
  \quad\text{and}\quad
  (\pi,0)
\]
modulo \(2\pi\). The downward equilibrium is elliptic; the upward equilibrium
is hyperbolic.

The separatrix is the energy level through the hyperbolic point:
\[
  H=mg\ell.
\]
On this level,
\[
  \frac{p_\theta^2}{2m\ell^2}
  =
  mg\ell(1+\cos\theta)
  =
  2mg\ell\cos^2\frac{\theta}{2}.
\]
Therefore
\[
  p_\theta
  =
  \pm2m\ell\sqrt{g\ell}\cos\frac{\theta}{2}
\]
for \(-\pi\leq\theta\leq\pi\). The separatrix divides libration orbits, which
oscillate around \(\theta=0\), from rotation orbits, which wind around the
circle. This picture is the one-degree-of-freedom ancestor of separatrix
splitting in perturbed Hamiltonian systems.

\section[Oscillator Actions In Polar Coordinates]{Action Variables For The
Isotropic Oscillator In Polar Coordinates}
\label{sec:detailed-oscillator-actions}

A canonical test case for action-angle coordinates is the two-dimensional
isotropic oscillator
\[
  H=
  \frac{p_r^2}{2m}
  +
  \frac{p_\theta^2}{2mr^2}
  +
  \frac12m\omega^2r^2.
\]
The angular momentum
\[
  L=p_\theta
\]
is conserved. The angular action is immediate:
\[
  I_\theta=\frac1{2\pi}\oint p_\theta\,\dd\theta=L.
\]
The radial action is
\[
  I_r
  =
  \frac1{\pi}
  \int_{r_-}^{r_+}
  \sqrt{
    2mE-\frac{L^2}{r^2}-m^2\omega^2r^2
  }\,\dd r.
\]
Rather than evaluate this integral by brute force, use a structural argument.
For the Cartesian oscillator,
\[
  H=\omega(I_x+I_y).
\]
The angular momentum generates rotations in the plane. For nonnegative \(L\),
the circular orbit has all its action in rotation and no radial oscillation.
On a circular orbit,
\[
  I_r=0,
  \qquad
  E=\omega L.
\]
Radial oscillation has twice the angular oscillator frequency: one full radial
cycle occurs while the Cartesian oscillator runs through half of its ellipse.
Therefore the Hamiltonian in polar actions must be
\[
  \boxed{
  H=\omega(2I_r+|I_\theta|).
  }
\]
Solving for the radial action,
\[
  \boxed{
  I_r=\frac{E}{2\omega}-\frac{|L|}{2}.
  }
\]
This formula is a useful reminder: action variables depend on the chosen basis
of cycles. Cartesian actions and polar actions are related by an integral
linear change of torus cycles, not by an arbitrary smooth coordinate change.

\section{Symmetric Free Top: Actions And Dense Motion}
\label{sec:detailed-symmetric-top-actions}

For a symmetric free rigid body with
\[
  I_1=I_2>I_3,
\]
use Euler angles \((\phi,\theta,\psi)\). The conserved quantities are
\[
  H,\qquad J^2,\qquad J_z.
\]
The momentum conjugate to \(\psi\) is the body-axis component of angular
momentum:
\[
  p_\psi=M_3.
\]
The momentum conjugate to \(\phi\) is the spatial vertical component:
\[
  p_\phi=J_z.
\]
For the symmetric top,
\[
  H
  =
  \frac{J^2-M_3^2}{2I_1}
  +
  \frac{M_3^2}{2I_3}.
\]
Solving for \(M_3^2\) gives
\[
  M_3^2
  =
  \frac{2I_1I_3}{I_1-I_3}
  \left(
    H-\frac{J^2}{2I_1}
  \right).
\]
Thus one action is
\[
  I_\psi=M_3
  =
  \sqrt{
  \frac{2I_1I_3}{I_1-I_3}
  \left(
    H-\frac{J^2}{2I_1}
  \right)
  },
\]
with sign chosen by the orientation of the \(\psi\)-cycle. Another action is
\[
  I_\phi=J_z.
\]
The remaining nutation action is
\[
  \boxed{
  I_\theta=J-I_\psi.
  }
\]
This compact result has a geometric explanation. The angular momentum vector
has fixed length \(J\). The body symmetry axis precesses around \(J\), and the
spin about the body axis contributes \(I_\psi\). The residual part of the
angular-momentum sphere swept by the nutation cycle is \(J-I_\psi\).

The angle equations are
\[
  \dot\alpha_j=\frac{\partial H}{\partial I_j}.
\]
Unless the three frequencies are rationally related, a trajectory is dense on
the corresponding invariant three-torus. The rigid body therefore gives a
physical example in which dense quasi-periodic motion occurs without chaos.

\section{Driven One-Angle Hamiltonian And Resonant Normal Form}
\label{sec:detailed-driven-resonance}

Consider the one-angle time-periodic perturbation model
\[
  H(I,\theta,t)
  =
  \frac12I^2+\epsilon\cos(n\theta-\nu t).
\]
The unperturbed frequency is
\[
  \dot\theta=I.
\]
Resonance occurs when the phase
\[
  \psi=n\theta-\nu t
\]
is slow, i.e.
\[
  nI-\nu\approx0.
\]
Thus
\[
  I_{\mathrm{res}}=\frac{\nu}{n}.
\]
To make the system autonomous, introduce the extended phase-space pair
\((t,E_t)\), with Hamiltonian
\[
  K=E_t+\frac12I^2+\epsilon\cos(n\theta-\nu t).
\]
Use the canonical slow angle
\[
  \psi=n\theta-\nu t
\]
and an action \(P\) measuring displacement from resonance. Locally set
\[
  I=\frac{\nu}{n}+P.
\]
Dropping constants and terms independent of the slow dynamics, the resonant
Hamiltonian has the pendulum form
\[
  \boxed{
  K_{\mathrm{res}}(P,\psi)
  =
  \frac12P^2+\epsilon\cos\psi.
  }
\]
Up to scale factors depending on the precise canonical normalization, this is
the universal local model of a first-order resonance. Its phase portrait has
elliptic islands, hyperbolic fixed points, and a separatrix. This is why
planetary mean-motion resonances, the standard map, and the driven pendulum all
share the same local geometry.

\section{Henon-Heiles Perturbation And The First Small Divisor}
\label{sec:detailed-henon-small-divisor}

The explicit Henon-Heiles calculation is carried out in
Section~\ref{sec:lab-henon-heiles}. The point needed here is the general
normal-form lesson: when the unperturbed oscillator has
\(\omega_0=(1,1)\), Fourier modes with \(k_1+k_2=0\) are resonant, so the
homological equation cannot remove them by division. They must be retained in
the resonant normal form, where they organize the islands and separatrices seen
in Poincare sections.

\section{Pressure As A Lagrange Multiplier In Incompressible Flow}
\label{sec:detailed-pressure-multiplier}

For incompressible Navier-Stokes,
\[
  \partial_tu+u\cdot\nabla u
  =
  -\frac1\rho\nabla p+\nu\Delta u,
  \qquad
  \nabla\cdot u=0.
\]
The pressure is not determined by a separate thermodynamic equation in the
incompressible model. It is the Lagrange multiplier enforcing
\(\nabla\cdot u=0\).

Take the divergence of the momentum equation. Since
\[
  \nabla\cdot\partial_tu=\partial_t(\nabla\cdot u)=0
  \quad\text{and}\quad
  \nabla\cdot\Delta u=\Delta(\nabla\cdot u)=0,
\]
we obtain the pressure Poisson equation
\[
  \boxed{
  -\frac1\rho\Delta p
  =
  \nabla\cdot(u\cdot\nabla u)
  =
  \partial_i u_j\,\partial_j u_i.
  }
\]
The last equality follows by expanding:
\[
  \partial_i(u_j\partial_j u_i)
  =
  (\partial_i u_j)(\partial_j u_i)
  +
  u_j\partial_j(\partial_i u_i),
\]
and the second term vanishes by incompressibility. Thus \(p\) is whatever
spatial field is needed to project the nonlinear acceleration back onto the
space of divergence-free vector fields.

\section{Orr-Sommerfeld Equation With Boundary Conditions}
\label{sec:detailed-orr-sommerfeld}

For a parallel base flow
\[
  U(y)e_x,
\]
linearized incompressible perturbations satisfy
\[
  \partial_t\tilde u
  +
  U\partial_x\tilde u
  +
  \tilde v U'e_x
  =
  -\nabla\tilde p+\frac1{\mathrm{Re}}\Delta\tilde u,
  \qquad
  \nabla\cdot\tilde u=0.
\]
For two-dimensional perturbations, introduce a streamfunction
\[
  \tilde u=\partial_y\psi,
  \qquad
  \tilde v=-\partial_x\psi.
\]
The perturbation vorticity is
\[
  \tilde\omega_z
  =
  \partial_x\tilde v-\partial_y\tilde u
  =
  -\Delta\psi.
\]
Taking the curl of the linearized momentum equation gives
\[
  \left(\partial_t+U\partial_x\right)\Delta\psi
  -
  U''\partial_x\psi
  =
  \frac1{\mathrm{Re}}\Delta^2\psi.
\]
Insert the normal mode
\[
  \psi(x,y,t)=\phi(y)e^{i\alpha x+\sigma t}.
\]
Then
\[
  \Delta\psi=(D^2-\alpha^2)\phi\,e^{i\alpha x+\sigma t},
\]
and the eigenvalue equation becomes
\[
  \boxed{
  (\sigma+i\alpha U)(D^2-\alpha^2)\phi
  -
  i\alpha U''\phi
  =
  \frac1{\mathrm{Re}}(D^2-\alpha^2)^2\phi.
  }
\]
No penetration means
\[
  \tilde v=-\partial_x\psi=0
  \quad\Rightarrow\quad
  \phi=0
\]
at a wall. No slip also requires
\[
  \tilde u=\partial_y\psi=0
  \quad\Rightarrow\quad
  D\phi=0.
\]
Thus the Orr-Sommerfeld equation is fourth order and needs the four wall
conditions \(\phi=D\phi=0\) at the two boundaries.

\section{Mechanical Connection: Curvature From A Small Shape Loop}
\label{sec:detailed-connection-curvature}

Let \(s^1,s^2\) be two shape coordinates and let the zero-momentum
reconstruction equation be
\[
  g^{-1}\dot g
  =
  -\mathcal A_\alpha(s)\dot s^\alpha.
\]
For a tiny rectangular loop based at \(s_0\), first move by
\(\Delta s^1\), then \(\Delta s^2\), then back by \(-\Delta s^1\), then back
by \(-\Delta s^2\). To second order in the side lengths, the net group element
is
\[
  g_{\mathrm{loop}}
  =
  \exp(-\mathcal A_2\Delta s^2)
  \exp(-\mathcal A_1\Delta s^1)
  \exp(\mathcal A_2\Delta s^2)
  \exp(\mathcal A_1\Delta s^1)
\]
plus the correction from the variation of \(\mathcal A_\alpha(s)\) across the
loop. The Baker-Campbell-Hausdorff formula gives the commutator contribution,
while the spatial variation gives derivative terms. The result is
\[
  g_{\mathrm{loop}}
  =
  \exp\left(
    -\mathcal F_{12}\Delta s^1\Delta s^2
    +
    O(|\Delta s|^3)
  \right),
\]
where
\[
  \boxed{
  \mathcal F_{12}
  =
  \partial_1\mathcal A_2-\partial_2\mathcal A_1
  +
  [\mathcal A_1,\mathcal A_2].
  }
\]
This is why curvature is the local density of geometric phase. Even if the
connection components are constant, a nonabelian structure group can give
nonzero holonomy through the commutator term. The helicopter rotor model in
Chapter~\ref{chap:worked-problem-labs} is exactly of this type.

\section{What These Examples Teach}
\label{sec:detailed-examples-summary}

\begin{itemize}
\item Central-force problems are solved by symmetry reduction before any
special orbit equation is guessed.
\item The Kepler problem is special because its conics and its
Laplace-Runge-Lenz vector are two expressions of the same hidden symmetry.
\item Separatrices are not just special curves in a picture; they organize
nearby motion and become the templates for chaotic transport after perturbation.
\item Action variables are geometric period integrals. They are useful because
they survive coordinate changes better than local phase coordinates do.
\item Small perturbations of integrable systems do not merely change numbers;
they change the geometry of phase space through resonances and separatrices.
\item Pressure in incompressible flow and multipliers in constrained mechanics
play the same structural role: they enforce constraints.
\item Gauge curvature measures the failure of local reconstruction rules to be
path independent.
\item A worked derivation is complete only when the final equation is tied back
to the original assumptions: which terms were neglected, which quantities are
held fixed, and which variables are observable.
\end{itemize}
